\documentclass[12pt,a4paper]{article}
\usepackage[utf8]{inputenc}
\usepackage[T1]{fontenc}
\usepackage{amsmath,amssymb}
\usepackage[margin=1in]{geometry}
\begin{document}
\title{Kuliah 11 - Shrinkage Methods dalam Sains Data}
\author{Yeni Herdiyeni}
\date{}
\maketitle

\maketitle
\tableofcontents
\newpage

\section{Kuliah 11 - Shrinkage Methods dalam Sains Data}

\textbf{Topik:} Ridge dan Lasso Regression  

\textbf{Mata Kuliah:} Pengantar Sains Data  

\textbf{Pengampu:} Yeni Herdiyeni, Program Studi Sarjana Kecerdasan Buatan, IPB University


\hrule


\subsection{1. Alur Pembahasan}

Materi ini dibahas secara bertahap melalui lima langkah utama:


\begin{enumerate}
\item Mengenal \textbf{OLS} sebagai model dasar regresi linear.
\end{enumerate}

\begin{enumerate}
\item Melihat permasalahan OLS pada data dengan fitur yang saling berkorelasi.
\end{enumerate}

\begin{enumerate}
\item Memahami bagaimana \textbf{Ridge} mengatasi masalah tersebut dengan menyusutkan semua koefisien.
\end{enumerate}

\begin{enumerate}
\item Memahami bagaimana \textbf{Lasso} menyusutkan koefisien sekaligus melakukan seleksi fitur.
\end{enumerate}


\hrule


\subsection{2. Use Case: Prediksi Hasil Panen Padi}

\subsubsection{2.1 Kasus}

Kita ingin memprediksi hasil panen padi berdasarkan sejumlah variabel lingkungan dan budidaya.


\textbf{Contoh fitur yang digunakan:}


\begin{tabular}{|l|l|}
\hline
Fitur & Keterangan \\
\hline
\texttt{rainfall} & Curah hujan \\
\hline
\texttt{humidity} & Kelembapan udara \\
\hline
\texttt{temperature} & Temperatur \\
\hline
\texttt{soil_moisture} & Kelembapan tanah \\
\hline
\texttt{irrigation} & Irigasi \\
\hline
\texttt{fertilizer_A} & Pupuk A \\
\hline
\texttt{fertilizer_B} & Pupuk B \\
\hline
\texttt{sunlight} & Cahaya matahari \\
\hline
\texttt{soil_ph} & pH tanah \\
\hline
\texttt{pest_index} & Indeks hama \\
\hline
\end{tabular}


\subsubsection{2.2 Tujuan}

Membangun model yang dapat memperkirakan hasil panen berdasarkan kombinasi seluruh variabel tersebut.


\subsubsection{2.3 Mengapa Kasus Ini Menarik?}

Faktor-faktor yang memengaruhi hasil panen biasanya tidak berdiri sendiri. Banyak variabel bergerak bersama dan saling berhubungan:



\textbf{Akibatnya:} Model menjadi lebih sulit membedakan mana pengaruh murni dari tiap variabel. Di sinilah persoalan \textbf{multikolinearitas} mulai muncul.


\hrule


\subsection{3. Apa Itu OLS?}

\subsubsection{3.1 Ordinary Least Squares (OLS)}

OLS adalah metode regresi linear yang mencari model terbaik dengan cara meminimalkan jumlah kuadrat error antara nilai aktual dan nilai prediksi:


\\[
\min_{\beta} \sum_{i=1}^{n} (y_i - \hat{y}_i)^2
\\]


Di mana:


\subsubsection{3.2 Model OLS pada Use Case}

\\[
\text{yield} = \beta_0 + \beta_1(\text{rainfall}) + \beta_2(\text{temperature}) + \dots + \beta_p(\text{pest index})
\\]


Koefisien $\beta_j$ menunjukkan seberapa besar kontribusi masing-masing fitur terhadap hasil panen.


\subsubsection{3.3 Cara Membaca OLS}


\subsubsection{3.4 Kekuatan OLS}


\hrule


\subsection{4. Permasalahan OLS: Multikolinearitas}

\subsubsection{4.1 Apa itu Multikolinearitas?}

\textbf{Multikolinearitas} adalah kondisi di mana beberapa prediktor (fitur) saling berhubungan secara linear. Pada data hasil panen, banyak prediktor saling berhubungan.


\subsubsection{4.2 Dampak Multikolinearitas pada OLS}

Koefisien OLS dapat berubah cukup besar ketika data sedikit berubah. Hal ini menunjukkan bahwa model memang bisa \emph{fit} pada data latih, tetapi belum tentu \textbf{stabil} untuk interpretasi dan prediksi baru.


\textbf{Makna visualisasi:} Setiap box dalam plot koefisien menunjukkan variasi koefisien dari banyak sampel ulang. Semakin lebar variasinya, semakin tidak stabil estimasi koefisien OLS.


\subsubsection{4.3 Mengapa Perlu Shrinkage?}

\textbf{Ide utamanya sederhana:}


\textbf{Shrinkage hadir untuk:}

\begin{enumerate}
\item Sekaligus menambahkan penalti pada besar koefisien.
\end{enumerate}

\begin{enumerate}
\item Membuat model lebih stabil dan lebih baik dalam generalisasi.
\end{enumerate}


\hrule


\subsection{5. Ridge Regression}

\subsubsection{5.1 Gagasan Utama Ridge}

Ridge menggunakan OLS sebagai dasar, lalu menambahkan penalti terhadap \textbf{kuadrat koefisien}. Dengan demikian, semua koefisien didorong untuk mengecil.


\subsubsection{5.2 Fungsi Objektif Ridge}

\\[
\min_{\beta_0, \beta} \sum_{i=1}^{n} \left( y_i - \beta_0 - \sum_{j=1}^{p} x_{ij}\beta_j \right)^2 + \lambda \sum_{j=1}^{p} \beta_j^2
\\]


Di mana:


\subsubsection{5.3 Pengaruh Nilai $\lambda$}

\begin{tabular}{|l|l|}
\hline
Nilai $\lambda$ & Efek \\
\hline
$\lambda = 0$ & Ridge sama seperti OLS \\
\hline
$\lambda$ membesar & Penalti makin kuat \\
\hline
$\lambda$ sangat besar & Semua koefisien mendekati nol, tetapi umumnya tidak tepat nol \\
\hline
\end{tabular}


\subsubsection{5.4 Intuisi Ridge pada Kasus Hasil Panen}

\textbf{Mengapa Ridge membantu?}


\begin{quote}
\textbf{Pesan sederhana:} Ridge berkata, "Semua fitur boleh tetap masuk, tetapi jangan ada koefisien yang terlalu ekstrem."\end{quote}


\subsubsection{5.5 Karakteristik Ridge}


\hrule


\subsection{6. Lasso Regression}

\subsubsection{6.1 Gagasan Utama Lasso}

Lasso juga berangkat dari OLS, tetapi penalti yang digunakan adalah \textbf{jumlah nilai absolut koefisien}. Penalti ini membuat sebagian koefisien dapat menjadi tepat nol.


\subsubsection{6.2 Fungsi Objektif Lasso}

\\[
\min_{\beta_0, \beta} \sum_{i=1}^{n} \left( y_i - \beta_0 - \sum_{j=1}^{p} x_{ij}\beta_j \right)^2 + \lambda \sum_{j=1}^{p} \lvert \beta_j \rvert
\\]


Di mana:


\subsubsection{6.3 Pengaruh Nilai $\lambda$}

\begin{tabular}{|l|l|}
\hline
Nilai $\lambda$ & Efek \\
\hline
$\lambda$ kecil & Hampir sama seperti OLS \\
\hline
$\lambda$ membesar & Beberapa koefisien menyusut kuat sampai nol \\
\hline
$\lambda$ sangat besar & Hanya sedikit fitur yang tersisa \\
\hline
\end{tabular}


\subsubsection{6.4 Intuisi Lasso pada Kasus Hasil Panen}

\textbf{Mengapa Lasso menarik?}


\begin{quote}
\textbf{Pesan sederhana:} Lasso berkata, "Jika suatu fitur tidak cukup penting, lebih baik koefisiennya dibuat nol."\end{quote}


\subsubsection{6.5 Efek Seleksi Fitur pada Lasso}

Semakin besar penalti $\lambda$, semakin banyak fitur yang dieliminasi dari model. Karena itu, Lasso sangat berguna saat kita juga ingin \textbf{menyederhanakan model}, bukan hanya meningkatkan kestabilan.


\hrule


\subsection{7. Perbandingan OLS, Ridge, dan Lasso}

\begin{tabular}{|l|l|l|l|}
\hline
Aspek & OLS & Ridge & Lasso \\
\hline
Fungsi tujuan & Minimalkan SSE & SSE + $\lambda \sum \beta_j^2$ & SSE + $\lambda \sum \lvert \beta_j \rvert$ \\
\hline
Penalti & Tidak ada & $L2$ & $L1$ \\
\hline
Koefisien & Tidak dibatasi & Menyusut, tapi tidak nol & Bisa menjadi tepat nol \\
\hline
Seleksi fitur & Tidak & Tidak & Ya \\
\hline
Multikolinearitas & Rentan & Tangguh & Tangguh \\
\hline
Kompleksitas model & Penuh & Penuh & Lebih ringkas \\
\hline
\end{tabular}


\subsubsection{Visualisasi Efek Penalti}


\hrule


\subsection{8. Kapan Menggunakan Ridge dan Kapan Lasso?}

\subsubsection{Gunakan Ridge Jika...}


\subsubsection{Gunakan Lasso Jika...}


\subsubsection{Kombinasi: Elastic Net}

Dalam praktik, sering digunakan \textbf{Elastic Net} yang menggabungkan penalti $L1$ dan $L2$:


\\[
\min_{\beta} \sum_{i=1}^{n} (y_i - \hat{y}_i)^2 + \lambda_1 \sum_{j=1}^{p} \lvert \beta_j \rvert  + \lambda_2 \sum_{j=1}^{p} \beta_j^2
\\]


Elastic Net memberikan keuntungan baik Ridge (stabil terhadap multikolinearitas) maupun Lasso (seleksi fitur).


\hrule


\subsection{9. Implementasi Python}

\subsubsection{9.1 Ridge Regression}

\begin{lstlisting}[language=Python]
from sklearn.linear_model import Ridge
from sklearn.preprocessing import StandardScaler
from sklearn.pipeline import make_pipeline

# Ridge dengan alpha = lambda
ridge_model = make_pipeline(
    StandardScaler(),
    Ridge(alpha=1.0)
)

ridge_model.fit(X_train, y_train)
y_pred_ridge = ridge_model.predict(X_test)
\end{lstlisting}

\subsubsection{9.2 Lasso Regression}

\begin{lstlisting}[language=Python]
from sklearn.linear_model import Lasso
from sklearn.preprocessing import StandardScaler
from sklearn.pipeline import make_pipeline

# Lasso dengan alpha = lambda
lasso_model = make_pipeline(
    StandardScaler(),
    Lasso(alpha=0.1)
)

lasso_model.fit(X_train, y_train)
y_pred_lasso = lasso_model.predict(X_test)
\end{lstlisting}

\begin{quote}
\textbf{Catatan:} Standardisasi fitur sangat penting sebelum menerapkan Ridge atau Lasso karena penalti sensitif terhadap skala fitur.\end{quote}


\hrule


\subsection{10. Interpretasi Hasil dalam Konteks Pertanian}

\subsubsection{10.1 Ridge: Semua Faktor Dipertimbangkan}

Dalam prediksi hasil panen, Ridge cocok ketika kita percaya bahwa banyak variabel lingkungan dan budidaya secara kolektif memengaruhi hasil panen, meskipun beberapa di antaranya saling berkorelasi.


\subsubsection{10.2 Lasso: Memilih Faktor Kunci}

Lasso cocok ketika kita ingin menyederhanakan model dan hanya menyimpan faktor-faktor yang paling berpengaruh. Misalnya, dari 10 fitur, mungkin hanya 4 atau 5 yang benar-benar dominan dalam memprediksi hasil panen.


\subsubsection{10.3 Kalimat Penutup}

\begin{quote}
Pada masalah prediksi hasil panen, \textbf{Ridge} membantu ketika banyak variabel sama-sama relevan, sedangkan \textbf{Lasso} membantu ketika kita juga ingin mengetahui fitur mana yang benar-benar penting.\end{quote}


\hrule


\subsection{11. Ringkasan}

\begin{enumerate}
\item \textbf{Multikolinearitas} dapat membuat koefisien OLS tidak stabil.
\end{enumerate}

\begin{enumerate}
\item \textbf{Ridge Regression} menambahkan penalti $L2$ untuk menyusutkan semua koefisien dan mengatasi multikolinearitas.
\end{enumerate}

\begin{enumerate}
\item \textbf{Lasso Regression} menambahkan penalti $L1$ untuk menyusutkan koefisien dan melakukan seleksi fitur otomatis.
\end{enumerate}

\begin{enumerate}
\item Gunakan \textbf{Ridge} jika banyak fitur sama-sama penting dan ada multikolinearitas.
\end{enumerate}

\begin{enumerate}
\item Gunakan \textbf{Lasso} jika ingin model yang lebih ringkas dan dapat diinterpretasikan.
\end{enumerate}

\begin{enumerate}
\item \textbf{Standardisasi fitur} wajib dilakukan sebelum menerapkan Ridge atau Lasso.
\end{enumerate}

\begin{enumerate}
\item \textbf{Elastic Net} menggabungkan keunggulan kedua metode.
\end{enumerate}



\end{document}
